% Section 6. Written by loop0009 item 09 (review section 3, "Section 5,
% Closing"). Every claim is from the body: Section 2 summary (rows), Section 3
% (communities), Sections 4.3-4.5 (rules), Section 5 (cost, values, certificates,
% dataset). The in-flight status is `python -m paper1.solver_fix_split
% --summary`, checked 2026-10-06. The price of the open classes is
% paper1/dataset.md section 6 (moved here from Section 5.4 by item 06).
\section{Conclusion}
\label{sec:closing}

We asked which equalities of the Linhares--Yanasse table are true, and what
follows when results, algorithms and benchmarks are moved between its twelve
problems through pathwidth. Eight rows are exact, counting the pathwidth row
itself: each of those eight problems is pathwidth, or pathwidth plus one, on
every input. Two rows, split bandwidth and edge search, are bands
one and two wide, and two, PLA folding and edge separation, are false as stated,
although a variant of each is exact (Table~\ref{tab:master}). Moving between
the problems paid twice. Read as a pathwidth algorithm, the best exact
open-stacks search turned out to rest on two false dominance theorems, and the
rule that corrects them was already known in the pathwidth literature. And one
certified pathwidth answers every problem in the exact core at once, which is
what made a single dataset for twelve problems possible, three of them
problems the table did not list, one from a fourth discipline, numerical
linear algebra.

\paragraph{What transfers, and what does not.} Within the exact core, anything
proved or computed for one problem holds for the others with the offset of
Table~\ref{tab:master}: an algorithm, a lower bound, a benchmark value. Across
a band only an approximation transfers. Split bandwidth and edge search stay
within one of the number of open stacks, as the table's literal reading says,
but not at a fixed offset, so an exact value of one does not give an exact
value of the other. Across a false row nothing transfers. Simple PLA folding
can need arbitrarily more tracks than the pathwidth suggests, and cutwidth is
not within a constant of pathwidth on stars. Section~\ref{sec:names} showed why
these distinctions went unchecked for two decades: the communities rarely read
one another, and most of what is known about the family is published under
``pathwidth''.

\paragraph{What the search teaches about dominance rules.} The two false
theorems were published in 2009 and restated in later work, and the
implementations that use them reported no failure. On every benchmark, and in
every run of our checks where the published rules lost the last solution at a
node, another branch found a solution (Section~\ref{sec:theorem}). It took a
search built for the purpose to find a whole instance with a wrong optimum,
34 customers and one stack too many (Figure~\ref{fig:wrong}). So the absence of
wrong answers on benchmarks was never evidence that the rules were sound. Three things in this paper bear on that. Reading the rules in pathwidth language put the correct rule, the
commitment lemma, beside the published one. Proving each rule on its own was
not enough: the order in which the rules run matters, because a cycle of true
links can still lose a node (Section~\ref{sec:theorem}). And a certificate
whose checker recomputes every premise rejects the false refutation that a
checker of the published premises accepts (Section~\ref{sec:certs}). The repair
is cheap, one bipartite matching per candidate and at most 2\% more nodes per
benchmark set, and it changed none of the 115
benchmark optima it re-proved.

\paragraph{What a reader can do differently.} Cite the corrected table,
Table~\ref{tab:master}, rather than the 2002 one. Look for results on MOSP and
gate matrix layout under ``pathwidth''. And in any implementation of the rules
of \citet{ChuStuckey2009}, including the code behind the optima of
\citet{Frinhani2018} that \citet{MartinYanassePinto2022} reuse
(Appendix~\ref{app:prior}), replace the published definite and better moves
with the matching test. The values those implementations report are not known
to be wrong, but their refutations are not proved.

\paragraph{Limits.} Four limits remain.
\begin{itemize}
\item The Lean model of the search and the C and Python code that run it are
  linked by testing and by certificates, not by proof
  (Section~\ref{sec:theorem}).
\item That an acyclic covering by checked steps is sound, which is the
  certificate checker's acceptance condition, is a paper argument
  (Section~\ref{sec:certs}).
\item LaPaugh's equality $\es=\pes$ is stated and not formalised. No row
  needs it (Section~\ref{sec:true}).
\item One reference of the table, \citet{LY2002ref5}, is not held. Its row,
  interval thickness, is proved from two other
  references of the table
  (Section~\ref{sec:true}).
\end{itemize}

\paragraph{Open problems.} One value of the MOSP corpus,
\path{Random-125-125-2-2_0}, has never been certified. As of 9 October 2026 the
split search is still refuting the value minus one and has found no order of
the smaller value; the other value that was open, \path{Random-125-125-2-3_0},
was certified on 9 October (Section~\ref{sec:practice}). Of the 17,714 classes of the dataset, 1,627 have
no certified pathwidth. One more budget step for every open class within the
size our implementation handles costs at most 500 core-hours, but on the Rome graphs each such step has closed only about
40\% of what it attempted, so the graph sets will not close by budget alone.
The 610 MOSP instances with 400 to 1,000 customers, and the 154 classes above
the implementation's limit of 1,024 vertices per component
(Section~\ref{sec:setup}), are beyond an exact search. Beyond budgets, a class is certified without search
only when a lower bound meets its witness (Section~\ref{sec:dataset}). The
open classes are therefore also a test set for stronger lower bounds on
pathwidth.
